\documentclass[11pt]{article}

\usepackage{amsmath, amssymb, graphicx, geometry}
\usepackage{booktabs}
\usepackage{listings}
\usepackage{xcolor}
\usepackage{hyperref}

\geometry{margin=1in}

\lstset{
  basicstyle=\ttfamily\small,
  keywordstyle=\color{blue},
  commentstyle=\color{gray},
  stringstyle=\color{red},
  breaklines=true,
  columns=fullflexible,
  keepspaces=true,
  language=C++
}

\newcommand{\file}[1]{\texttt{#1}}
\newcommand{\type}[1]{\texttt{#1}}
\newcommand{\fn}[1]{\texttt{#1}}

\title{Validation Infrastructure for Environment-Responsive\\
Coarse-Grained Bead Dynamics}
\author{VSEPR-SIM Project}
\date{\today}

\begin{document}

\maketitle

\begin{abstract}
The companion documents establish the anisotropic surface-mapped bead
(ASB) and the environment-responsive extension (ERB), introducing an
extended bead state $\mathbb{B}_B = \{B,\, X_B\}$ with fast observables
$(\rho_B, C_B, P_{2,B})$ and a slow internal state variable $\eta_B$.
This document specifies the \emph{testing and validation infrastructure}
developed to verify the environment-responsive layer and its downstream
analysis pipeline: synthetic scene construction, trajectory-based time
evolution, behavioural assertion frameworks, Monte Carlo robustness
harnesses, structured formation studies at scale $N > 200$, ensemble
proxy validation, macro property precursor channel testing, property
learning pipeline verification, and property calibration and target
legitimacy validation.  The infrastructure is not auxiliary to the
model---it is a co-equal component of model development.  Key findings
include: (i) the seven reported state variables reduce to four
independent degrees of freedom, (ii)~large-$N$ runs ($N = 64, 125, 216$)
are numerically stable and bounded, (iii)~edge versus bulk differentiation
emerges naturally from the environment descriptors, (iv)~the causal
chain geometry $\to$ observables $\to$ target $\to$ $\eta$ is now
empirically confirmed, and (v)~the full pipeline from per-bead state
through ensemble proxies, precursor channels, property learning, and
calibrated prediction is validated across 1013 tests.
\end{abstract}

\tableofcontents
\newpage

% ============================================================================
\section{Introduction and Position in the Document Series}
% ============================================================================

This is the fourth document in the coarse-grained methodology series:

\begin{center}
\begin{tabular}{cll}
    \textbf{Code} & \textbf{Title} & \textbf{Establishes} \\
    \midrule
    CGL & Coarse-Grained and Multi-Scale Layer
        & Architecture, scale transitions, ensembles,
          ensemble proxies, precursors, pipeline, calibration \\
    ASB & Anisotropic Surface-Mapped Beads
        & Descriptor, orientation, interaction engine \\
    ERB & Environment-Responsive Bead Dynamics
        & $X_B$, fast observables, $\eta_B$, kernel modulation \\
    \textbf{TIF} & \textbf{Testing Infrastructure (this document)}
        & \textbf{Validation methodology, empirical findings} \\
\end{tabular}
\end{center}

\subsection{Why a Separate Document?}

Environment-responsive dynamics introduces a new category of model risk.
The bead state now includes a dynamical variable ($\eta_B$) that evolves
via a relaxation equation, couples back into the interaction kernels, and
is expected to exhibit memory, hysteresis, and path dependence.  None of
these properties can be verified by unit-testing individual functions in
isolation.  They require:

\begin{enumerate}
    \item \textbf{Synthetic scene construction:} controlled spatial
          configurations of beads with known geometry, spacing, and
          orientation distributions
    \item \textbf{Trajectory-based verification:} time-evolution of $\eta$
          under static or controlled conditions, with full observable
          recording
    \item \textbf{Behavioural assertions:} monotonicity, boundedness,
          convergence, and sensitivity checks that go beyond equality tests
    \item \textbf{Statistical robustness:} thousands of randomised trials
          verifying that the model does not blow up under stochastic input
    \item \textbf{Large-$N$ stability:} structured runs with $N > 200$
          beads demonstrating that collective behaviour is well-posed
\end{enumerate}

Each of these requires infrastructure that is itself a scientific
contribution.  This document specifies that infrastructure, catalogues the
results it has produced, and establishes the validation gates that must be
passed before proceeding to higher-level studies.

\subsection{Testing Philosophy}

The testing infrastructure follows a \textbf{state-first, not
chemistry-first} philosophy.  Rather than constructing atomistic molecular
systems, mapping them to beads, and then testing the environment layer, the
infrastructure manufactures bead scenes directly---bypassing atomistic
chemistry entirely.

This design choice is deliberate:

\begin{itemize}
    \item It isolates the environment-responsive layer from mapping errors
    \item It enables deterministic, reproducible test configurations
    \item It allows systematic exploration of neighbourhood geometries
          that may not correspond to any real molecule
    \item It makes tests fast (microseconds, not minutes)
\end{itemize}

The philosophy is not that atomistic tests are unimportant, but that
\emph{model-level} validation must be independent of the pipeline that
feeds it.

% ============================================================================
\section{Scene Construction Methodology}
\label{sec:scenes}
% ============================================================================

A \emph{scene} is a collection of synthetic beads, each carrying a
position $\mathbf{R}_B$, an orientation axis $\hat{n}_B$, and an initial
internal state $\eta_B^{(0)}$.  Scenes are the fundamental test objects.

\subsection{The \type{SceneBead} Representation}

Each synthetic bead is a lightweight struct:
\begin{equation}
    \text{SceneBead} = \bigl\{
        \mathbf{R},\;
        \hat{n},\;
        \texttt{has\_orientation},\;
        \eta^{(0)}
    \bigr\}
\end{equation}

This is intentionally minimal.  The bead carries no mass, no charge, no
descriptor coefficients, and no channel decomposition.  It carries only
what the environment-responsive layer needs to compute: positions,
orientations, and initial state.

\subsection{Deterministic Scene Builders}

The following canonical scenes provide controlled geometric configurations
with exact, reproducible coordinates:

\begin{center}
\begin{tabular}{lcp{6.5cm}}
    \textbf{Scene} & \textbf{$N$} & \textbf{Geometry} \\
    \midrule
    Isolated       & 1   & Single bead at origin \\
    Pair           & 2   & Two beads along $x$-axis at separation $d$ \\
    Linear stack   & $n$ & $n$ beads along $z$-axis, uniform spacing \\
    T-shape        & 3   & Center + aligned ($z$) + perpendicular ($x$) \\
    Square         & 4   & Four beads at corners of square in $xy$-plane \\
    Triangle       & 3   & Equilateral triangle in $xy$-plane \\
    Tetrahedron    & 4   & Regular tetrahedron vertices \\
    Ring           & $n$ & $n$ beads on circle in $xy$-plane \\
    Dense shell    & $n{+}1$ & Fibonacci spiral on sphere + central bead \\
    Cubic lattice  & $n^3$ & Perfect cubic grid \\
    Line bundle    & $n$ & Centered $z$-column, all aligned \\
    Layered slab   & $n^2 k$ & $k$ square planes stacked along $z$ \\
\end{tabular}
\end{center}

Each builder is a pure function: given geometric parameters, it returns a
deterministic scene.  No random state, no seed, no external dependency.

\subsection{Randomised Scene Generators}

For robustness testing, randomised generators produce controlled
distributions of bead positions and orientations.  All use a single
deterministic PRNG (\type{Xorshift32}) seeded explicitly, ensuring exact
reproducibility.

\begin{center}
\begin{tabular}{lp{8cm}}
    \textbf{Generator} & \textbf{Distribution} \\
    \midrule
    Random cluster    & $N$ beads uniformly in a box, random orientations \\
    Random cloud      & Center at origin + $N{-}1$ scattered in box \\
    Random shell      & Center at origin + $N$ on spherical shell \\
    Separated cloud   & Rejection-sampled: minimum inter-bead distance guaranteed \\
    Biased stack      & Column along $z$ with controlled jitter and alignment bias \\
    Perturbed lattice & Perfect lattice + bounded random displacement \\
    MC scene          & \type{MCSceneConfig}: $n \in [n_{\min}, n_{\max}]$
                        neighbours in $[r_{\min}, r_{\max}]$ shell \\
\end{tabular}
\end{center}

\subsection{Scene Transforms}

Two rigid transforms are provided for testing invariance properties:

\paragraph{Translation.}
$\mathbf{R}_B \to \mathbf{R}_B + \Delta\mathbf{R}$ for all beads.
Orientations are unaffected.  Environment observables must be
translation-invariant, so the environment state of a translated scene must
be identical to the original.

\paragraph{Rotation.}
Both positions and orientations are rotated via Rodrigues' formula about
an axis $\hat{k}$ by angle $\theta$:
\begin{equation}
    \mathbf{v}' = \mathbf{v}\cos\theta
    + (\hat{k} \times \mathbf{v})\sin\theta
    + \hat{k}(\hat{k} \cdot \mathbf{v})(1 - \cos\theta)
\end{equation}

Applied to both $\mathbf{R}_B$ and $\hat{n}_B$.  This is a rigid rotation:
inter-bead distances and relative orientations are preserved.  All
environment observables must be rotation-invariant.

\subsection{Scene Validation}

Before running any test, a scene can be validated for structural sanity:

\begin{center}
\begin{tabular}{lp{8cm}}
    \textbf{Check} & \textbf{What it catches} \\
    \midrule
    Minimum pair distance    & Catastrophic overlaps, degenerate geometry \\
    Duplicate positions      & Copy errors, bad generation \\
    Unit normal verification & Non-unit orientations after rotation or perturbation \\
    Centroid                 & Drift, asymmetric generation \\
    Bounding box             & Scale errors, unexpected extent \\
\end{tabular}
\end{center}

These checks prevent debugging generated scenes by staring at coordinates.

\subsection{Neighbour List Construction}

The neighbour list builder returns \emph{all} other beads in the scene as
neighbours for a given bead, regardless of distance.  The cutoff filtering
is performed inside the environment state kernel
(\fn{compute\_fast\_observables}), which applies distance-based Gaussian
weighting controlled by $r_c$ and $\sigma_\rho$.

This design ensures that scene construction is completely independent of
kernel parameters.  The same scene can be evaluated under different cutoff
regimes without regeneration.

% ============================================================================
\section{Variable Hierarchy}
\label{sec:hierarchy}
% ============================================================================

The environment-responsive bead reports seven state variables.  A central
finding of the testing infrastructure is that these seven variables are
\emph{not} seven independent degrees of freedom.

\subsection{The Seven Reported Variables}

\begin{center}
\begin{tabular}{llll}
    \textbf{Variable} & \textbf{Symbol} & \textbf{Type} & \textbf{Range} \\
    \midrule
    Local density           & $\rho$        & Fast observable    & $[0, \infty)$ \\
    Normalised density      & $\hat{\rho}$  & Derived            & $[0, 1]$ \\
    Coordination number     & $C$           & Fast observable    & $[0, N{-}1]$ \\
    Orientational order     & $P_2$         & Fast observable    & $[-\tfrac{1}{2}, 1]$ \\
    Normalised order        & $\hat{P}_2$   & Derived            & $[0, 1]$ \\
    Target function         & $f$           & Composite          & $[0, 1]$ \\
    Internal state          & $\eta$        & Dynamical variable & $[0, 1]$ \\
\end{tabular}
\end{center}

\subsection{Independence Structure}

Systematic sweeps across canonical scenes (Suite~\#3, Phase~1) reveal:

\begin{itemize}
    \item $\hat{\rho}$ is a deterministic rescaling of $\rho$:
          $\hat{\rho} = \rho / \rho_{\max}$
    \item $\hat{P}_2$ is a deterministic rescaling of $P_2$:
          $\hat{P}_2 = (P_2 + \tfrac{1}{2}) / \tfrac{3}{2}$
    \item $f$ is a linear combination of normalised forms:
          $f = \alpha\hat{\rho} + \beta\hat{P}_2$
\end{itemize}

Therefore the system has four independent variables:
\begin{equation}
\boxed{
    \text{Independent:}\quad \rho,\; C,\; P_2,\; \eta
    \qquad
    \text{Derived:}\quad \hat{\rho} = g(\rho),\;
    \hat{P}_2 = h(P_2),\;
    f = \alpha\, g(\rho) + \beta\, h(P_2)
}
\end{equation}

This is not a deficiency---it is a clarification.  The normalised forms
and the target function are not independent scientific degrees of freedom.
They are intermediate computational quantities.

\subsection{Causal Chain}

The empirically confirmed causal structure of the environment-responsive
layer is:

\begin{equation}
    \underbrace{
        \text{geometry, spacing, orientation}
    }_{\text{scene configuration}}
    \;\longrightarrow\;
    \underbrace{
        \rho,\; C,\; P_2
    }_{\text{measured local state}}
    \;\longrightarrow\;
    \underbrace{
        \hat{\rho},\; \hat{P}_2
    }_{\text{normalised forms}}
    \;\longrightarrow\;
    \underbrace{
        f
    }_{\text{target}}
    \;\longrightarrow\;
    \underbrace{
        \eta
    }_{\text{memory}}
\end{equation}

This chain means that $\eta$ is not a primary descriptor of the scene.  It
is a \emph{stateful memory variable} that integrates the scene's geometric
properties through the intermediate observable layer.  If someone asks
``what determines $\eta$?'', the answer is not ``the scene directly'' but
``the scene determines descriptors, descriptors determine the target, the
target drives $\eta$.''

This causal interpretation should govern all subsequent test design,
parameter fitting, and scientific reporting.

% ============================================================================
\section{Test Suite Architecture}
\label{sec:suites}
% ============================================================================

The validation infrastructure is organised into five test suites, each
targeting a different aspect of model behaviour.  All suites share a common
assertion pattern and report format.

\subsection{Suite Overview}

\begin{center}
\begin{tabular}{lrlp{5.5cm}}
    \textbf{Suite} & \textbf{Tests} & \textbf{File} & \textbf{Focus} \\
    \midrule
    Suite \#1  & 60  & \file{test\_anisotropic\_model.cpp}
        & Channel kernels, per-$\ell$ decomposition, SH rotation \\
    Suite \#2  & 77  & \file{test\_environment\_suite2.cpp}
        & Observable correctness, $\eta$ dynamics, constitutive
          response, coupled pipeline, stress/pathology \\
    Track 2    & 39  & \file{test\_track2\_robustness.cpp}
        & Monte Carlo robustness across 4000+ random scenes \\
    Suite \#3  & 52  & \file{test\_formation\_suite3.cpp}
        & Structured $N > 10$ formation studies, atlas, large-$N$ \\
    Suite \#4  & 113 & \file{test\_formation\_regimes\_suite4.cpp}
        & Dynamic regimes, invariance, condition-to-structure \\
    Suite \#5  & 131 & \file{test\_ensemble\_proxy\_suite5.cpp}
        & Emergent effective medium mapping, ensemble response proxies,
          directed proxy validation, delta/validity/reference \\
    Suite \#6  & 85  & \file{test\_macro\_precursor\_suite6.cpp}
        & Macro property precursor channels, directed perturbation,
          confidence model, inter-channel consistency \\
    Suite \#7  & 109 & \file{test\_property\_pipeline\_suite7.cpp}
        & Property learning pipeline, dataset construction,
          feature vectorization, regression, synthetic supervision \\
    Suite \#8  & 193 & \file{test\_property\_calibration\_suite8.cpp}
        & Property calibration, target legitimacy, confidence
          decomposition, OOD detection, promotion/demotion \\
    CLI Suite  & 72  & \file{test\_cg\_cli.cpp}
        & Integration tests for all CG console commands \\
    Bead Layers & 82 & \file{test\_bead\_layers.cpp}
        & Visualization data model, bead layer rendering,
          inspector overlay, detail view modes \\
    \midrule
    \textbf{Total} & \textbf{1013} & & \\
\end{tabular}
\end{center}

All 1013 tests pass.  No test is skipped or expected to fail.

\subsection{Suite \#2: Environment Observable Correctness}

Suite~\#2 is the most structurally detailed, organised into five groups:

\begin{description}
    \item[Group A: Environment observable correctness.]
        Neighbourhood sensing, translation invariance, rotation invariance,
        distance sweeps, isolated bead special cases.

    \item[Group B: Internal state dynamics.]
        $\eta$ evolution under constant environment, $\tau$-dependence,
        determinism under identical initial conditions, timestep
        convergence.

    \item[Group C: Constitutive response.]
        Kernel modulation factors, backward compatibility (all $\gamma_k
        = 0$ recovers unmodulated kernels), modifier boundedness.

    \item[Group D: Coupled pipeline.]
        Multi-step evolution, history-dependent behaviour, emergent
        ordering under sustained alignment.

    \item[Group E: Stress and pathology.]
        Extreme values ($\eta < 0$, $\eta > 1$), NaN/Inf guards, empty
        neighbour lists, very large distances, degenerate orientations.
\end{description}

\subsection{Track 2: Monte Carlo Robustness}

Track~2 runs large ensembles of randomised local bead scenes and verifies
statistical properties.  It is a property-based statistical validation,
not a unit test.

\begin{center}
\begin{tabular}{rl}
    \textbf{Ensemble} & \textbf{$N$ trials} \\
    \midrule
    Isotropic random    & 1000 \\
    Aligned dense       & 1000 \\
    Mixed distribution  & 2000 \\
\end{tabular}
\end{center}

Properties verified across all 4000 trials:
\begin{itemize}
    \item $\eta \in [0, 1]$ at all timesteps
    \item No NaN or Inf in any observable
    \item Isotropic environments do not spuriously activate ordered state
    \item Aligned dense environments correctly activate ordering
    \item Kernel modulations remain finite and sign-correct
    \item Higher alignment bias produces higher $P_2$ (monotonic trend)
\end{itemize}

\subsection{Suite \#3: Structured Formation Studies}

Suite~\#3 consists of three phases:

\paragraph{Phase 1: Single-variable response atlas.}
For each of seven environment variables, sweep a controlling scene
parameter (spacing, from 2.0 to 15.0) while holding all others fixed.
Run across six canonical scenes.  Record output range, monotonicity, and
perturbation sensitivity.

\paragraph{Phase 2a: Pairwise coupling maps.}
Sweep two variables simultaneously on coarse $5 \times 5$ grids.  Four
coupling maps: $(\rho \times \eta)$, $(C \times P_2)$,
$(P_2 \times \hat{P}_2)$, $(\rho \times f)$.  Detect nonlinear
transitions, saturation, and divergence.

\paragraph{Phase 2b: Large-$N$ runs.}
$N = 64$, $125$, $216$ beads with controlled initial conditions across six
scene types: perfect lattice, perturbed lattice, line bundle, layered slab,
shell initialisation, random cloud.  All 18 runs stable and bounded.

% ============================================================================
\section{Behavioural Assertion Framework}
\label{sec:assertions}
% ============================================================================

Traditional unit tests assert equality: \emph{expected output equals actual
output}.  Environment-responsive testing requires \emph{behavioural}
assertions that characterise qualitative properties of trajectories and
distributions.

\subsection{Trajectory Assertions}

\begin{center}
\begin{tabular}{lp{7cm}}
    \textbf{Assertion} & \textbf{Meaning} \\
    \midrule
    \fn{is\_monotone\_increasing}
        & Sequence is non-decreasing (within $10^{-15}$ tolerance) \\
    \fn{is\_monotone\_decreasing}
        & Sequence is non-increasing \\
    \fn{is\_bounded}$(v, a, b)$
        & All values in $[a, b]$ \\
    \fn{approaches\_value}$(v, t, \epsilon)$
        & Final value within $\epsilon$ of target $t$ \\
    \fn{has\_variation}
        & Not all values are identical (system responded) \\
    \fn{is\_finite}
        & No NaN or Inf anywhere in the sequence \\
\end{tabular}
\end{center}

These are not ad hoc checks.  Each corresponds to a physical property of
the relaxation equation:
\begin{itemize}
    \item Monotone increasing $\eta$ when $f > \eta$ (driving toward
          target from below)
    \item Bounded in $[0, 1]$ by construction of the clamped relaxation
    \item Convergence to $f$ at steady state
    \item Variation confirms the system is not degenerate
    \item Finiteness guards against numerical pathology
\end{itemize}

\subsection{Statistical Assertions}

For ensemble-level verification:

\begin{center}
\begin{tabular}{lp{7cm}}
    \textbf{Tool} & \textbf{Purpose} \\
    \midrule
    \fn{compute\_stats}
        & Mean, variance, percentiles (5th, 50th, 95th), NaN/Inf counts \\
    \fn{mean\_increases\_with}$(x, y)$
        & Binned monotonic trend: mean of $y$ increases with sorted $x$ \\
    \fn{groups\_separated}$(A, B, \delta)$
        & Mean of group $B$ exceeds mean of group $A$ by at least $\delta$ \\
\end{tabular}
\end{center}

These tools enable the Track~2 robustness harness to make statistical
claims across thousands of trials without hand-inspecting individual runs.

\subsection{Sweep Harness}

The direct sweep harness automates single-variable parameter studies:

\begin{enumerate}
    \item Accept a scene-generating function $S(v)$ parameterised by
          sweep variable $v$
    \item Linearly sample $v$ from $v_{\min}$ to $v_{\max}$
    \item For each sample: build scene, evolve all beads for $n$ steps,
          record bead~0 state
    \item Return a vector of \type{ResponseRecord} with all seven
          environment variables
\end{enumerate}

This enables systematic parameter studies with a single function call,
replacing ad hoc loops in test files.

% ============================================================================
\section{Trajectory Runner}
\label{sec:runner}
% ============================================================================

The trajectory runner is the central execution primitive.  It time-evolves
a single bead within a static (frozen-position) scene and records the full
history of all observables.

\subsection{Single-Bead Trajectory}

Given a scene, a bead index, and initial $\eta^{(0)}$:

\begin{enumerate}
    \item Build the neighbour list for the target bead (all other beads
          in the scene, unfiltered)
    \item For each timestep $t = 0, 1, \ldots, n{-}1$:
    \begin{enumerate}
        \item Compute fast observables $(\rho, C, P_2)$ from the current
              neighbour configuration
        \item Normalise: $\hat{\rho} = \rho / \rho_{\max}$,
              $\hat{P}_2 = (P_2 + \tfrac{1}{2}) / \tfrac{3}{2}$
        \item Compute target: $f = \alpha\hat{\rho} + \beta\hat{P}_2$
        \item Integrate: $\eta(t{+}1) = f + (\eta(t) - f)\,
              e^{-\Delta t / \tau}$
        \item Record: $(\eta, f, \rho, P_2, C)$ appended to trajectory
    \end{enumerate}
    \item Return the complete \type{EtaTrajectory}
\end{enumerate}

Because positions are frozen, the fast observables are constant across
timesteps.  Only $\eta$ evolves.  This makes the trajectory runner a test
of the relaxation dynamics in isolation, cleanly separated from the
question of whether positions would change under the modified forces.

\subsection{Multi-Bead Pass}

For large-$N$ studies, the multi-bead runner evolves all beads
simultaneously.  At each timestep, every bead's environment state is
updated from its own neighbour list.  The initial $\eta$ for each bead is
read from the \type{SceneBead::eta} field, enabling non-uniform
initialisation for hysteresis and memory tests.

% ============================================================================
\section{Key Empirical Findings}
\label{sec:findings}
% ============================================================================

The testing infrastructure has produced several findings that inform
subsequent model development.

\subsection{Variable Redundancy}

As documented in Section~\ref{sec:hierarchy}, the seven reported state
variables contain only four independent degrees of freedom.  The
normalised variables $\hat{\rho}$ and $\hat{P}_2$ are deterministic
rescalings of $\rho$ and $P_2$ respectively, and the target function $f$
is a linear combination of the normalised forms.

This finding has two consequences:
\begin{enumerate}
    \item Subsequent analysis should treat $\rho$, $C$, $P_2$, and $\eta$
          as the primary state, not all seven variables
    \item If the target function is later generalised to a nonlinear form,
          the redundancy structure will change and must be re-evaluated
\end{enumerate}

\subsection{Large-$N$ Stability Gate}

Eighteen structured runs at $N = 64$, $125$, and $216$ across six scene
types (perfect lattice, perturbed lattice, line bundle, layered slab,
shell, random cloud) all produced:

\begin{itemize}
    \item $\eta \in [0, 1]$ for all beads at all timesteps
    \item No NaN or Inf in any observable
    \item Bounded, finite output across all configurations
\end{itemize}

This clears the large-$N$ stability gate.  The environment-responsive
layer does not exhibit numerical instability under any tested structured
initialisation at scale $N \leq 216$.

\subsection{Edge vs.\ Bulk Differentiation}

On the $6 \times 6 \times 6 = 216$ cubic lattice:

\begin{center}
\begin{tabular}{lcc}
    \textbf{Property} & \textbf{Bulk} & \textbf{Edge} \\
    \midrule
    $\rho$  & high  & low \\
    $\eta$  & high  & low \\
    $C$     & high  & low \\
\end{tabular}
\end{center}

This differentiation is physically expected: interior beads have more
neighbours, higher density, and therefore higher $\eta$ at steady state.
Edge beads have fewer neighbours and lower observables.

The fact that this emerges naturally from the environment descriptors---
without explicit surface detection or boundary tagging---validates that the
observables carry physically meaningful topological information.

\subsection{Coordination Histogram Structure}

On a $5 \times 5 \times 5 = 125$ cubic lattice, the coordination number
$C$ exhibits a structured distribution corresponding to geometric position:

\begin{center}
\begin{tabular}{ll}
    \textbf{Position} & \textbf{Coordination} \\
    \midrule
    Corner beads    & lowest $C$ \\
    Edge beads      & intermediate $C$ \\
    Face beads      & higher $C$ \\
    Interior beads  & highest $C$ \\
\end{tabular}
\end{center}

The histogram shows at least two distinct coordination values, confirming
that the environment descriptors resolve topological position within a
structured formation.

\subsection{Monte Carlo Robustness}

Across 4000 randomised trials:
\begin{itemize}
    \item No numerical divergence in any trial
    \item Isotropic random environments produce low $\eta$ (no spurious
          activation)
    \item Aligned dense environments correctly produce high $\eta$
    \item Higher alignment bias monotonically increases $P_2$
    \item Dense shells produce higher $C$ than sparse shells
\end{itemize}

These are not individual test results---they are statistical properties
of the model verified across an ensemble.

% ============================================================================
\section{Infrastructure Architecture}
\label{sec:arch}
% ============================================================================

\subsection{Header Organisation}

The testing infrastructure is split into three headers:

\begin{center}
\begin{tabular}{lp{8cm}}
    \textbf{Header} & \textbf{Contents} \\
    \midrule
    \file{scene\_builders.hpp}
        & \type{SceneBead}, \type{Xorshift32}, all scene builders,
          \fn{translate\_scene}, \fn{rotate\_scene}, scene validation
          helpers, \fn{build\_neighbours}, \fn{set\_initial\_eta} \\
    \file{test\_runners.hpp}
        & \type{EtaTrajectory}, \fn{run\_trajectory},
          \fn{run\_all\_beads}, behavioural assertions, \type{StatSummary},
          \fn{compute\_stats}, \type{RunRecord},
          \fn{run\_robustness\_trial}, \type{ResponseRecord},
          \fn{run\_single\_variable\_sweep},
          \type{CoordinationHistogram} \\
    \file{scene\_factory.hpp}
        & Backward-compatible aggregator: includes both of the above \\
\end{tabular}
\end{center}

The split separates concerns:
\begin{itemize}
    \item \emph{What} to test (scene construction) is independent of
          \emph{how} to test (trajectory execution, assertion checking)
    \item New scene types can be added without touching the runner
    \item New assertion helpers can be added without touching the builders
    \item Existing consumers include \file{scene\_factory.hpp} and
          automatically receive both
\end{itemize}

\subsection{PRNG Policy}

All randomised generation uses a single PRNG implementation
(\type{Xorshift32}) with explicit seeding.  This ensures:

\begin{itemize}
    \item Exact reproducibility: same seed $\Rightarrow$ same scene
    \item Platform independence: no reliance on \texttt{std::mt19937}
          distribution implementation
    \item Debuggability: any failing trial can be reproduced by seed
    \item Determinism: required by the project's anti-black-box philosophy
\end{itemize}

\subsection{Dependency Structure}

\begin{center}
\begin{tabular}{lcl}
    \textbf{Header} & $\longrightarrow$ & \textbf{Depends on} \\
    \midrule
    \file{scene\_builders.hpp}
        & $\longrightarrow$
        & \file{environment\_state.hpp}, \file{state.hpp} (Vec3) \\
    \file{test\_runners.hpp}
        & $\longrightarrow$
        & \file{scene\_builders.hpp}, \file{environment\_coupling.hpp} \\
    \file{scene\_factory.hpp}
        & $\longrightarrow$
        & Both of the above \\
\end{tabular}
\end{center}

The dependency on \file{environment\_state.hpp} is necessary because
\fn{build\_neighbours} returns \type{NeighbourInfo} structs, and the
runners call \fn{update\_environment\_state}.  The dependency on
\file{environment\_coupling.hpp} is needed only by
\fn{run\_robustness\_trial} for kernel modulation reports.

% ============================================================================
\section{Validation Gates}
\label{sec:gates}
% ============================================================================

The testing infrastructure defines explicit gates that must be passed
before proceeding to higher-level studies.

\subsection{Gate 1: Observable Correctness (Suite \#2, Group A)}

All fast observables must be:
\begin{itemize}
    \item Translation-invariant
    \item Rotation-invariant (when both positions and orientations are
          rotated)
    \item Monotonic in expected directions (density increases with
          neighbour count)
    \item Zero for isolated beads ($\rho = 0$, $C = 0$, $P_2$
          undefined/zero)
    \item Finite for all non-degenerate inputs
\end{itemize}

\textbf{Status: Passed.} 77/77 tests in Suite~\#2.

\subsection{Gate 2: Dynamical Stability (Suite \#2, Groups B--D)}

The slow variable $\eta$ must:
\begin{itemize}
    \item Remain in $[0, 1]$ at all times
    \item Converge to $f$ at steady state
    \item Exhibit monotone approach when starting below/above target
    \item Be deterministic under identical initial conditions
    \item Be insensitive to timestep within reasonable range
\end{itemize}

\textbf{Status: Passed.} All dynamical tests in Suite~\#2 Groups B--D.

\subsection{Gate 3: Robustness Under Stochastic Input (Track 2)}

The model must not produce:
\begin{itemize}
    \item NaN or Inf for any random scene configuration
    \item $\eta$ outside $[0, 1]$
    \item Spurious ordering in isotropic environments
    \item Failure to order in aligned dense environments
\end{itemize}

\textbf{Status: Passed.} 39/39 tests across 4000 random trials.

\subsection{Gate 4: Large-$N$ Stability (Suite \#3, Phase 2b)}

At $N = 64, 125, 216$:
\begin{itemize}
    \item All beads bounded
    \item All observables finite
    \item Edge/bulk differentiation present
    \item Coordination histogram structured
\end{itemize}

\textbf{Status: Passed.} 18/18 large-$N$ runs stable.

\subsection{Gate 5: Variable Structure (Suite \#3, Phase 1)}

The response atlas must confirm:
\begin{itemize}
    \item Identified redundancy pattern ($\hat{\rho}$, $\hat{P}_2$, $f$
          are derived)
    \item Four independent degrees of freedom
    \item Monotonic or bounded response across all scenes
    \item Finite output across the full sweep range
\end{itemize}

\textbf{Status: Passed.} 52/52 tests in Suite~\#3.

\subsection{Suite \#4: Dynamic Formation Regimes}

Suite~\#4 (113 tests) implements the dynamic, invariance, and regime-mapping
studies described below.  It is organised into four phases:

\paragraph{Phase 4A: Dynamic validation (21 tests).}
Hysteresis under spacing ramps (compress then decompress), alignment ramps
(fully aligned vs.\ random orientation), high-$\eta$ vs.\ low-$\eta$
initialisation comparison, relaxation time extraction, and $\tau$-variation
effects.  Key findings:
\begin{itemize}
    \item Compressed and decompressed scenes remain bounded and finite
    \item Low-$\eta$ and high-$\eta$ initialisations converge to the same
          final mean $\eta$ (difference $< 0.05$), confirming the system is
          a first-order relaxation without multistability
    \item Relaxation timing $t_{10\%} \le t_{5\%} \le t_{1\%}$ with all
          thresholds reached within the run window
    \item Faster $\tau$ produces faster relaxation, as expected
\end{itemize}

\paragraph{Phase 4B: Invariance validation (18 tests).}
Translation invariance, rigid rotation covariance (Rodrigues'),
bead-order permutation invariance (Fisher--Yates shuffle), seeded
reproducibility, and rotation composition.  All passed:
\begin{itemize}
    \item $\eta$, $\rho$, $C$, $P_2$ are exactly invariant under
          translation
    \item All observables and $\eta$ are invariant under rigid rotation
          (tolerance $10^{-10}$)
    \item Sorted per-bead observables are identical under permutation
    \item Fixed seeds produce bit-identical scenes and trajectories;
          different seeds produce different scenes
    \item Two composed $90^\circ$ rotations match a single $180^\circ$
          rotation
\end{itemize}

\paragraph{Phase 4C: Formation-regime atlas (29 tests).}
Sweeps across six scene families (lattice, perturbed lattice, line bundle,
layered slab, shell, separated cloud), five $\alpha/\beta$ weight settings,
five $\tau$ values, and five alignment bias levels.  Bulk/edge
classification via coordination threshold.  Results:
\begin{itemize}
    \item All spacing sweeps bounded, finite, and responsive
    \item $\alpha/\beta$ weight variation affects $\eta$
    \item All $\tau$ sweep runs stable
    \item Alignment bias measurably changes $\eta$ and $P_2$
    \item Non-degenerate geometries show clear bulk/edge differentiation
          with bulk $C \ge$ edge $C$
    \item Uniform-distance geometries (shell) correctly show degenerate
          coordination
\end{itemize}

\paragraph{Phase 4D: Medium-to-large-$N$ scaling (45 tests).}
$N = 64$, $125$, $216$ across four scene types (perfect lattice, perturbed
lattice, shell, separated cloud).  Plus cross-$N$ scaling consistency and
formation history convergence:
\begin{itemize}
    \item All 12 large-$N$ runs bounded and finite
    \item Bulk interior $\eta$ consistent across $N$ (max difference $<
          0.15$)
    \item Late-stage $\eta$ variance $\le$ early-stage variance,
          confirming convergence
    \item Coordination histograms valid at all sizes
\end{itemize}

% ============================================================================
\section{Completed Extensions (Stage 4)}
\label{sec:extensions}
% ============================================================================

The following were planned extensions that have now been implemented and
verified in Suite~\#4.

\subsection{Formation-Under-Conditions Experiments}

Systematic variation of:
\begin{itemize}
    \item Density and spacing (compression/expansion)
    \item Alignment bias (ordered $\leftrightarrow$ disordered)
    \item Shell count and local coordination
    \item Target function weights ($\alpha$, $\beta$)
    \item Relaxation timescale ($\tau$)
    \item Initial $\eta$ distributions
\end{itemize}

\textbf{Status: Completed.} All sweeps bounded, finite, and responsive.

\subsection{Hysteresis and Memory Tests}

Tested path dependence:
\begin{itemize}
    \item Compress then decompress spacing
    \item Aligned vs.\ random orientations
    \item Initialise $\eta$ low vs.\ high under same scene
\end{itemize}

\textbf{Status: Completed.} The system converges to the same final state
regardless of initial $\eta$ (no multistability detected), confirming
pure first-order relaxation dynamics.

\subsection{Time-to-Equilibrium Quantification}

Measured:
\begin{itemize}
    \item Time to within 1\%, 5\%, 10\% of final value
    \item Correct ordering: $t_{10\%} \le t_{5\%} \le t_{1\%}$
    \item $\tau$-dependent relaxation speed
    \item Convergence of late-stage variance (monotone relaxation)
\end{itemize}

\textbf{Status: Completed.} No oscillation, overshoot, or numerical fragility
detected across all tested regimes.

\subsection{Invariance Validation}

Newly established invariance proofs:
\begin{itemize}
    \item Translation invariance (exact)
    \item Rigid rotation covariance (tolerance $10^{-10}$)
    \item Bead-order permutation invariance (sorted observables)
    \item Seeded reproducibility (bit-identical)
\end{itemize}

\textbf{Status: Completed.} All invariances confirmed.

% ============================================================================
\section{Emergent Effective Medium Mapping (Stage 5)}
\label{sec:eem}
% ============================================================================

The ensemble proxy layer bridges per-bead environment state to
macroscopic response characterisation.  It does not claim to predict
continuum-scale constitutive properties; it provides\emph{proxy fields}
computed from the distributions and spatial organisation of bead-state
variables.

\subsection{Architectural Position}

\begin{center}
\begin{tabular}{l}
    Environment-state evaluation (per-bead) \\
    $\downarrow$ \\
    Ensemble statistics / spatial field summaries \\
    $\downarrow$ \\
    Macroscopic response proxies \\
    $\downarrow$ \\
    Later constitutive or transport models \\
\end{tabular}
\end{center}

The analysis module lives at \file{coarse\_grain/analysis/ensemble\_proxy.hpp}
and produces an \texttt{EnsembleProxySummary} from a vector of
\texttt{EnvironmentState} and bead positions.

\subsection{Summary Sections}

The \texttt{EnsembleProxySummary} struct contains seven sections:

\begin{enumerate}[label=(\Alph*)]
    \item \textbf{First moments:} mean $\rho$, $\hat{\rho}$, $C$,
          $P_2$, $\hat{P}_2$, $\eta$, $f$, and mean $|\eta - f|$
    \item \textbf{Second moments:} var($\rho$), var($C$), var($P_2$),
          var($\eta$), var($f$)
    \item \textbf{Edge/bulk contrasts:} bulk and edge bead counts,
          gap in $\rho$, $C$, $P_2$, $\eta$
    \item \textbf{Occupancy fractions:} fraction of beads above
          coordination, $\eta$, and alignment thresholds
    \item \textbf{Spatial autocorrelation:} neighbour-pair
          autocorrelation of $\eta$, $\rho$, $P_2$
    \item \textbf{Response proxies:} five scalar proxies bounded
          in $[0, 1]$
    \item \textbf{Diagnostics:} all-finite and all-bounded flags
\end{enumerate}

\subsection{Response Proxy Definitions}

Five macroscopic response proxies are defined:

\begin{description}
    \item[Cohesion] $(\bar{\hat{\rho}} + \bar{\eta}
        + (1 - \overline{|\eta - f|})) / 3$.
        Higher for dense, stabilised, low-mismatch systems.
    \item[Uniformity] $1 - (d_\eta + d_\rho) / 2$ where
        $d = \mathrm{clamp}(2\sqrt{\mathrm{var}}, 0, 1)$.
        Higher for homogeneous systems.
    \item[Texture] $\overline{\hat{P}_2}$.
        Higher for oriented, anisotropic organisation.
    \item[Stabilization] $\bar{\eta}\,(1 - \overline{|\eta - f|})$.
        Higher for converged, well-adapted states.
    \item[Surface sensitivity] Mean fractional $|$bulk $-$ edge$|$
        gaps in $\rho$, $C$, $\eta$.
        Higher when bulk and edge environments are sharply distinct.
\end{description}

\subsection{Suite \#5: Validation (131 tests)}

Suite~\#5 validates the ensemble proxy layer across seven phases:

\begin{center}
\begin{tabular}{llr}
    \textbf{Phase} & \textbf{Focus} & \textbf{Tests} \\
    \midrule
    5A & Computation correctness (means, variances, bounds) & 30 \\
    5B & Physical sense (ordering, response, coherence) & 14 \\
    5C & Large-$N$ stability ($N = 64, 125, 216$) & 32 \\
    5D & Invariance (translation, rotation, permutation, seed) & 16 \\
    5E & Sensitivity (sweeps, families, parameter response) & 10 \\
    5F & Directed proxy validation (one perturbation per proxy) & 10 \\
    5G & Delta, validity, reference (convergence, thresholds, normalisation) & 19 \\
    \midrule
    \textbf{Total} & & \textbf{131} \\
\end{tabular}
\end{center}

Key verified properties:
\begin{itemize}
    \item All five proxies remain in $[0, 1]$ across all tested scenes
          and parameters
    \item Dense lattices produce higher cohesion than sparse ones
    \item Aligned stacks produce higher texture than random clouds
    \item Converged systems have higher stabilisation than fresh ones
    \item All proxies are translation, rotation, and permutation invariant
    \item Results are seeded-reproducible (bit-identical)
    \item Large-$N$ proxies remain finite, bounded, and consistent
          across $N = 64 \to 216$
\end{itemize}
\label{sec:relation}
% ============================================================================

\begin{center}
\begin{tabular}{ll}
    \textbf{Companion document} & \textbf{This document} \\
    \midrule
    ERB defines $X_B = \{\rho, C, P_2, \eta\}$
        & Verifies $X_B$ correctness, stability, invariance \\
    ERB specifies relaxation dynamics
        & Tests convergence, boundedness, monotonicity \\
    ERB predicts emergent behaviour
        & Observes edge/bulk differentiation, coordination structure \\
    ERB introduces six new parameters
        & Sweeps parameters, confirms no blow-up \\
    ERB claims hysteresis and metastability
        & Suite~\#4 confirms relaxation, no multistability \\
    ERB + TIF establish variable hierarchy
        & Suite~\#5 constructs ensemble proxies from those variables \\
    CGL defines precursor channels
        & Suite~\#6 validates directional correctness and confidence \\
    CGL defines property pipeline
        & Suite~\#7 validates regression, ranking, and split discipline \\
    CGL defines calibration layer
        & Suite~\#8 validates confidence decomposition and legitimacy \\
\end{tabular}
\end{center}

The relationship is explicit: ERB and CGL make theoretical claims, TIF
provides the machinery to test them, and the test results either
confirm the claims or reveal where the theory requires amendment.

% ============================================================================
\section{Suite \#6: Macro Property Precursor Channels (Stage 6)}
\label{sec:suite6}
% ============================================================================

Suite~\#6 validates the macro property precursor layer
(\file{coarse\_grain/analysis/macro\_precursor.hpp}), which converts
ensemble-level proxies into eight property-like precursor channels.

\subsection{Precursor Channel Architecture}

The eight channels---rigidity-like, ductility-like, brittleness-like,
cohesion integrity-like, thermal transport-like, electrical transport-like,
surface reactivity-like, and fracture susceptibility-like---are intermediate
latent descriptors grouped by material-behaviour hypothesis.  All carry the
\texttt{-like} suffix to indicate they are evidence variables, not physical
properties.

Intermediate helpers (interface penalty, normalised correlation length,
adaptation capacity, texture lock, anisotropy index) are computed and
exposed for traceability before any channel formula is evaluated.

\subsection{Test Organisation}

\begin{center}
\begin{tabular}{llr}
    \textbf{Phase} & \textbf{Focus} & \textbf{Tests} \\
    \midrule
    6A & Channel computation, bounds, empty input & 22 \\
    6B & Directed perturbation (one test per channel) & 16 \\
    6C & Inter-channel consistency and relationships & 12 \\
    6D & Confidence model, degradation, provenance & 15 \\
    6E & Large-$N$ stability and cross-$N$ consistency & 20 \\
    \midrule
    \textbf{Total} & & \textbf{85} \\
\end{tabular}
\end{center}

All 85 tests pass.

% ============================================================================
\section{Suite \#7: Property Learning Pipeline (Stage 7)}
\label{sec:suite7}
% ============================================================================

Suite~\#7 validates the property learning pipeline
(\file{coarse\_grain/analysis/property\_pipeline.hpp}), which maps proxy
distributions and precursor channels to property-scale outputs with
explicit uncertainty quantification.

\subsection{Three-Layer Separation}

The pipeline enforces strict separation:
\begin{description}
    \item[Descriptor layer] Proxy summaries and precursor channels from
        simulation data.  No learned parameters.
    \item[Calibration layer] Model fitting from exported datasets.
        No simulation logic.
    \item[Inference layer] Prediction from fitted models.  No fitting
        and no simulation.
\end{description}

Formula fusion is prohibited.

\subsection{Test Organisation}

\begin{center}
\begin{tabular}{llr}
    \textbf{Phase} & \textbf{Focus} & \textbf{Tests} \\
    \midrule
    7A & Dataset row validity, targets, provenance & 15 \\
    7B & Vectorization determinism, ablation, missing policy & 15 \\
    7C & Linear fit, coefficients, prediction, abstention, attribution & 15 \\
    7D & RMSE/MAE, Spearman, Kendall, grouped splits & 15 \\
    7E & Ordinal labels, monotone sweeps, end-to-end pipeline & 49 \\
    \midrule
    \textbf{Total} & & \textbf{109} \\
\end{tabular}
\end{center}

All 109 tests pass.

% ============================================================================
\section{Suite \#8: Property Calibration and Target Legitimacy (Stage 8)}
\label{sec:suite8}
% ============================================================================

Suite~\#8 validates the property calibration and target legitimacy layer
(\file{coarse\_grain/analysis/property\_calibration.hpp}), which provides
a disciplined framework for calibrated property prediction with explicit
confidence decomposition and promotion/demotion policy.

\subsection{Calibration Discipline}

The layer introduces:
\begin{itemize}
    \item Five first-wave property families (rigidity, ductility,
          brittleness, thermal transport, electrical transport)
    \item Target-definition contracts with sign expectations and domain
          constraints
    \item Calibration profiles for interpretable output transformation
    \item Three supervision regimes (contract-supervised,
          synthetic-calibrated, externally calibrated)
    \item Five-component confidence decomposition (input, precursor,
          coverage, attribution, calibration)
    \item Out-of-distribution detection via feature envelopes
    \item Property-specific synthetic curricula
    \item Formal promotion/demotion policy for the \texttt{-like} suffix
\end{itemize}

\subsection{Test Organisation}

\begin{center}
\begin{tabular}{llr}
    \textbf{Phase} & \textbf{Focus} & \textbf{Tests} \\
    \midrule
    8A & Property family selection, naming, indexing & 16 \\
    8B & Target contracts, sign expectations, domain & 17 \\
    8C & Calibration profiles, transform, clamping & 12 \\
    8D & Supervision regime taxonomy & 4 \\
    8E & Confidence decomposition, status, attribution & 23 \\
    8F & OOD detection, envelope, extrapolation & 14 \\
    8G & Synthetic curricula, monotonicity, regime tagging & 14 \\
    8H & Legitimacy states, promotion, demotion, splits & 93 \\
    \midrule
    \textbf{Total} & & \textbf{193} \\
\end{tabular}
\end{center}

All 193 tests pass.

% ============================================================================
\section{Conclusion}
% ============================================================================

The testing infrastructure for environment-responsive coarse-grained bead
dynamics is not a peripheral convenience---it is a co-equal scientific
instrument.  It has:

\begin{enumerate}
    \item Established 1013 tests across eleven suites, all passing
    \item Identified the variable hierarchy: four independent degrees of
          freedom ($\rho$, $C$, $P_2$, $\eta$), three derived quantities
          ($\hat{\rho}$, $\hat{P}_2$, $f$)
    \item Confirmed the causal chain: geometry $\to$ observables $\to$
          target $\to$ $\eta$
    \item Cleared the large-$N$ stability gate at $N = 216$
    \item Demonstrated edge/bulk differentiation and coordination
          structure
    \item Verified robustness across 4000+ randomised trials
    \item Confirmed translation, rotation, and permutation invariance
    \item Completed formation-under-conditions, hysteresis, relaxation
          timing, and regime-atlas experiments (Suite~\#4)
    \item Established that the system is a pure first-order relaxation
          with no multistability or path-dependent memory
    \item Constructed and validated five macroscopic response proxies
          (cohesion, uniformity, texture, stabilization, surface
          sensitivity) from bead-state ensemble statistics (Suite~\#5,
          131 tests including directed validation and convergence diagnostics)
    \item Validated eight macro property precursor channels with directed
          perturbation tests and confidence modelling (Suite~\#6, 85 tests)
    \item Built and validated a property learning pipeline with dataset
          construction, 31-feature vectorization, ridge regression,
          grouped splits, and synthetic supervision (Suite~\#7, 109 tests)
    \item Established a property calibration and target legitimacy
          framework with five-component confidence decomposition,
          out-of-distribution detection, and formal promotion/demotion
          policy governing the \texttt{-like} suffix (Suite~\#8, 193 tests)
    \item Provided a complete, validated pipeline from per-bead environment
          state through ensemble proxies, precursor channels, property
          learning, and calibrated prediction
\end{enumerate}

The model is not merely defined.  It is tested, measured, and
characterised.  The testing infrastructure ensures that subsequent
development---including constitutive property determination, external
calibration, and discovery workflows---is grounded in empirical evidence
rather than theoretical assumption.

% ====================================================================
\section{Automated Visualization Infrastructure}
\label{sec:automated-viz}
% ====================================================================

The visualization checkpoint system (\file{test\_viz.hpp}) was extended
from a manual, blocking viewer into an automation-capable test
visualization infrastructure.  The original system---in which
\fn{CGVizViewer::run()} opens a GLFW/ImGui window and blocks until the
user manually closes it---is now augmented with three capabilities:

\begin{enumerate}
    \item \textbf{Timed display.}
          A configurable \fn{timeout\_seconds} field in \type{VizConfig}
          causes the viewer window to auto-close after the specified
          duration.  The render loop checks elapsed time each frame via
          \fn{glfwGetTime()} and issues \fn{glfwSetWindowShouldClose()}
          when the timeout expires.  Manual interaction (orbit, zoom,
          overlay cycling) remains fully active during the timed
          interval.

    \item \textbf{Command queue.}
          A sequence of \type{VizCommand} structs may be attached to the
          \type{VizConfig}.  During each iteration of the render loop,
          the viewer consumes commands from the front of the queue:
          \begin{itemize}
              \item Non-wait commands (\fn{SetOverlay}, \fn{SelectBead},
                    \fn{OrbitCamera}, \fn{SetAxes}, etc.) execute
                    immediately and are batched—multiple consecutive
                    non-wait commands all apply within a single frame.
              \item \fn{Wait} commands pause queue processing for the
                    specified number of seconds.  The viewer continues
                    rendering and accepting manual input during the
                    wait.
              \item \fn{Close} terminates the window programmatically.
          \end{itemize}
          After all commands are consumed, the viewer remains open until
          timeout or manual close, enabling a ``script then inspect''
          workflow.

    \item \textbf{Fluent builder (\type{VizSequence}).}
          A fluent builder class in \file{test\_viz.hpp} provides a
          human-readable API for constructing command sequences:
          \begin{lstlisting}
test_viz::VizSequence seq;
seq.overlay(test_viz::overlay::density)
   .wait(2.0f)
   .orbit(0.5f, 0.0f)
   .overlay(test_viz::overlay::memory)
   .select(3)
   .detail(true)
   .wait(1.5f)
   .close();
VIZ_AUTOMATED("Sweep test", state, seq);
          \end{lstlisting}
          The builder compiles to a no-op stub when
          \texttt{BUILD\_VIS=OFF}, preserving API surface for headless
          builds.
\end{enumerate}

\subsection{Supported Commands}

Table~\ref{tab:viz-commands} lists the full set of
\type{VizCommand} operations.

\begin{table}[h]
\centering
\caption{Automation command types for the viewer command queue.}
\label{tab:viz-commands}
\begin{tabular}{@{}lp{7cm}@{}}
\toprule
\textbf{Command} & \textbf{Effect} \\
\midrule
\fn{SetOverlay}      & Switch scalar overlay (none, $\rho$, $C$, $P_2$, $\eta$, $f^*$) \\
\fn{SetAxes}         & Toggle orientation axis rendering \\
\fn{SetNeighbours}   & Toggle neighbour edge rendering \\
\fn{SelectBead}      & Select bead by index ($-1$ to deselect) \\
\fn{SetDetailMode}   & Enter/exit multi-layer detail inspection \\
\fn{SetViewMode}     & Change detail view mode (Shell, Scaffold, Cutaway, \ldots) \\
\fn{OrbitCamera}     & Rotate camera by $(\Delta\theta, \Delta\phi)$ \\
\fn{ZoomCamera}      & Zoom camera by a multiplicative factor \\
\fn{PanCamera}       & Pan camera target by $(\Delta x, \Delta y)$ \\
\fn{ResetCamera}     & Re-center and auto-distance on bead centroid \\
\fn{Wait}            & Pause command processing for $t$ seconds \\
\fn{Close}           & Close the viewer window \\
\bottomrule
\end{tabular}
\end{table}

\subsection{Macro Summary}

The following macros are provided for test code:

\begin{table}[h]
\centering
\caption{Test visualization macros (all are no-ops when
  \texttt{VSEPR\_TEST\_VIZ} is not set).}
\label{tab:viz-macros}
\begin{tabular}{@{}lp{8cm}@{}}
\toprule
\textbf{Macro} & \textbf{Purpose} \\
\midrule
\fn{VIZ\_CHECKPOINT}          & Manual display of a \type{CGSystemState} \\
\fn{VIZ\_CHECKPOINT\_OVERLAY} & Manual display with scalar overlay \\
\fn{VIZ\_SCENE}               & Manual display of a \type{SceneBead} vector \\
\fn{VIZ\_SCENE\_OVERLAY}      & Manual display of scene with overlay \\
\fn{VIZ\_TIMED}               & Timed auto-close display \\
\fn{VIZ\_TIMED\_OVERLAY}      & Timed display with overlay \\
\fn{VIZ\_AUTOMATED}           & Scripted command sequence \\
\fn{VIZ\_AUTOMATED\_TIMEOUT}  & Scripted sequence with safety timeout \\
\fn{VIZ\_SCENE\_TIMED}        & Scene timed display \\
\fn{VIZ\_SCENE\_AUTOMATED}    & Scene with scripted sequence \\
\bottomrule
\end{tabular}
\end{table}

\subsection{Architecture Position}

The automation layer sits between test code and the GLFW render loop:

\begin{center}
\begin{tabular}{c}
Test code \\
$\downarrow$ \\
\fn{VIZ\_AUTOMATED} / \fn{VIZ\_TIMED} macros \\
$\downarrow$ \\
\fn{test\_viz::show\_automated()} / \fn{show\_timed()} \\
$\downarrow$ \\
\fn{VizConfig} with \fn{commands} + \fn{timeout\_seconds} \\
$\downarrow$ \\
\fn{CGVizViewer::run()} (automation-aware main loop) \\
$\downarrow$ \\
GLFW/ImGui window with live rendering
\end{tabular}
\end{center}

Manual input is never disabled.  Users can interact freely during
automated sequences, and the automation status (remaining timeout,
current command index, active wait progress) is displayed in the
ImGui controls panel.

\end{document}
